% Informe final de la actividad evaluable — compilar con `make informe`
% (Tectonic; descarga los paquetes LaTeX necesarios en la primera ejecución).
\documentclass[10pt,twocolumn]{article}

\usepackage[a4paper,top=2.2cm,bottom=2.4cm,left=1.9cm,right=1.9cm,columnsep=0.75cm]{geometry}
\usepackage[utf8]{inputenc}
\usepackage[T1]{fontenc}
\usepackage[spanish,es-tabla,es-noquoting]{babel}
\usepackage{amssymb}
\usepackage{newtxtext,newtxmath}
\usepackage{microtype}
\usepackage{booktabs}
\usepackage{graphicx}
\usepackage{xcolor}
\usepackage{titlesec}
\usepackage{caption}
\usepackage{fancyhdr}
\usepackage{enumitem}
\usepackage{balance}
\usepackage[hidelinks]{hyperref}

\definecolor{tinta}{RGB}{28,36,52}
\definecolor{acento}{RGB}{31,78,121}

\titleformat{\section}{\large\bfseries\color{acento}}{\thesection}{0.6em}{}
\titleformat{\subsection}{\normalsize\bfseries}{\thesubsection}{0.6em}{}
\titlespacing{\section}{0pt}{1.1em}{0.5em}
\titlespacing{\subsection}{0pt}{0.9em}{0.35em}

\captionsetup{font=small,labelfont={bf,color=acento},skip=4pt}
\setlist{nosep,leftmargin=1.15em}

\pagestyle{fancy}
\fancyhf{}
\fancyhead[L]{\small\itshape Aurum Market: búsqueda semántica y control de catálogo}
\fancyhead[R]{\small\thepage}
\renewcommand{\headrulewidth}{0.4pt}

\newcommand{\code}[1]{\texttt{\small #1}}

\title{\bfseries\color{tinta} Aurum Market: búsqueda semántica\\ y control de catálogo}
\author{Manuel Pérez\\
\small Evaluación · Bases de Datos Vectoriales\\
\small \today}
\date{}

\begin{document}

\twocolumn[
  \begin{@twocolumnfalse}
    \maketitle
    \vspace{-2.2em}
    \begin{center}
    \begin{minipage}{0.86\textwidth}
    \small
    \textbf{Resumen} --- El sistema entregado convierte los 15.000 productos de
    Aurum Market en un servicio de descubrimiento medible: recupera por
    intención, filtra por marca dentro de la propia consulta a la base de
    datos, absorbe altas, actualizaciones y bajas sin duplicar registros, y
    detecta altas potencialmente duplicadas con una regla calibrada solo con
    datos de desarrollo. El motor elegido es \textbf{Qdrant} (local, Docker)
    con su SDK nativo; la representación ganadora es
    \textbf{multilingual-e5-base} codificando \textbf{título + marca + color}
    en lugar del campo \code{text} completo --- decisiones que salieron de una
    rejilla de experimentos (composición, objetivo de entrenamiento y capacidad,
    con BM25 como baseline y Gemini Embedding 2 preparado como opción API), no
    del catálogo de modelos. Cada cifra de este
    informe se regenera desde \code{config/run\_config.yaml} y queda
    justificada, con tablas y gráficos, en la serie de notebooks de I+D
    \code{actividad\_00}--\code{05}.
    \end{minipage}
    \end{center}
    \vspace{1.6em}
  \end{@twocolumnfalse}
]

\section{Contrato del sistema y baseline}

El sistema expone una interfaz única (\code{DiscoveryService} y la CLI
\code{make search}) que recibe una consulta y devuelve resultados
normalizados: \code{product\_id}, posición, título, marca, \code{score}
nativo y su semántica (\code{similarity}, mayor es mejor). Los dos recorridos
de negocio comparten esa interfaz:

\begin{itemize}
  \item \textbf{Descubrimiento}: top-$k$ configurable, con filtro opcional de
  marca que viaja dentro de la consulta a la base de datos.
  \item \textbf{Control de altas}: la ficha entrante se codifica como
  documento, la base vectorial genera los candidatos y una regla de umbral
  decide si existe duplicado, señalando el \code{product\_id} concreto.
\end{itemize}

Como referencia interpretable se implementó un \textbf{baseline léxico BM25}
\cite{robertson2009} ($k_1{=}1.5$, $b{=}0.75$, tokenización en minúsculas y
sin tildes) sobre el mismo corpus. El baseline no es decorativo: fija el
listón que la representación densa debe superar y, como se ve en la
sección~\ref{sec:representacion}, no lo supera en todas las métricas, lo que
delimita honestamente dónde aporta la semántica.

La figura~\ref{fig:arquitectura} resume la arquitectura: datos
$\rightarrow$ representación (embeddings con manifiesto SHA-256)
$\rightarrow$ Qdrant (colección persistente con HNSW y payload indexado)
$\rightarrow$ servicios de búsqueda y duplicados $\rightarrow$ evaluación y
artefactos.

\begin{figure*}[t]
  \centering
  \includegraphics[width=0.92\textwidth]{arquitectura.pdf}
  \caption{Arquitectura del motor de descubrimiento y control de catálogo.}
  \label{fig:arquitectura}
\end{figure*}

\section{Representación vectorial}
\label{sec:representacion}

\textbf{Texto codificado.} Cada producto se representa con una composición
documentada y saneada: los valores vacíos son información ausente (nunca la
cadena \code{"nan"}) y los campos vacíos se omiten. Se compararon dos
composiciones y tres modelos locales \cite{wang2024,reimers2019}, más el
baseline léxico, siempre con búsqueda exacta para aislar la representación del
índice (tabla~\ref{tab:experimentos}); una configuración adicional con Gemini
Embedding 2 (API de la sesión 01) queda preparada y se construye solo si
existe \code{GEMINI\_API\_KEY}. El enunciado no obliga a modelos locales
(``no existe ventaja por utilizar cloud frente a local''); el final local es
una decisión de diseño: evaluación sin coste para el corrector y
regeneración completa sin clave ni binarios grandes en el repositorio.

\begin{table}[h]
  \centering\small
  \caption{Experimentos de representación sobre las 8 consultas de
  desarrollo (búsqueda exacta). En negrita, la configuración elegida.}
  \label{tab:experimentos}
  \begin{tabular}{@{}lccc@{}}
    \toprule
    Experimento & nDCG@10 & Recall@10 & MRR@10 \\
    \midrule
    bm25\_full\_text & 0.558 & 0.187 & 0.646 \\
    e5\_small\_full & 0.533 & 0.205 & 0.588 \\
    e5\_small\_title & 0.555 & 0.223 & 0.698 \\
    \textbf{e5\_base\_title} & \textbf{0.563} & \textbf{0.248} & \textbf{0.792} \\
    minilm\_full & 0.273 & 0.095 & 0.310 \\
    \bottomrule
  \end{tabular}
\end{table}

Tres conclusiones sostienen la elección de \code{e5\_base\_title}:

\begin{enumerate}
  \item \textbf{El campo \code{text} completo perjudica a E5.} Las fichas
  arrastran \emph{keyword-stuffing} (títulos repetidos, listas de variantes,
  HTML residual). Codificar solo título+marca+color mejora las tres métricas
  frente al texto completo ($+0.02$ nDCG, $+0.018$ recall, $+0.11$ MRR). La
  suciedad no se corrigió a mano: se neutralizó eligiendo qué codificar, que
  es una decisión trazable del experimento.
  \item \textbf{Cambiar de modelo se analizó, no solo se citó.} MiniLM, sin
  prefijos y entrenado para paráfrasis genérica, se hunde (nDCG 0.273):
  confirma que el prefijo \code{query:}/\code{passage:} y el entrenamiento de
  recuperación de E5 son la parte que importa, no la dimensión (ambos son
  384d).
  \item \textbf{La capacidad paga, una vez arreglada la composición.} Sobre el
  mismo texto ganador, e5-base (768d) añade $+0.008$ nDCG, $+0.025$ recall y
  $+0.094$ MRR frente a e5-small, con coste local trivial (6~s de embeddings
  con GPU, p50 de búsqueda ${\sim}2.7$ frente a ${\sim}2.3$~ms).
\end{enumerate}

\textbf{Métrica, normalización y score.} Los embeddings se L2-normalizan al
codificar y la colección usa distancia \textbf{coseno}; con vectores
unitarios el score devuelto por Qdrant es la similitud coseno en $[-1,1]$ y
``mayor es mejor''. Esa semántica se conserva de extremo a extremo
(\code{score\_kind="similarity"} en el contrato de resultados) y nunca se
compara con distancias.

Los embeddings se persisten en \code{data/embeddings/} con un manifiesto que
fija modelo, prefijos, composición, dimensión y checksums SHA-256 de entradas
y salidas, de modo que cualquier matriz puede auditarse contra el CSV del que
salió.

\section{Índice y base de datos}
\label{sec:indice}

\textbf{Esquema.} Una colección Qdrant \cite{qdrant}
(\code{aurum-market-eval-catalogo}) con un vector sin nombre de 768d,
distancia coseno, e ID de punto igual al \code{record\_id} UUIDv5 del
catálogo ---el mismo UUID que el manifiesto de datos declara, verificado en
carga---. El payload conserva \code{product\_id}, \code{title},
\code{brand}, \code{color}, \code{locale}, \code{catalog\_version} y
\code{active}; \code{brand} tiene índice de payload \emph{keyword} creado
junto con la colección para que el filtrado sea parte del plan de búsqueda.
Los valores ausentes se guardan como cadena vacía, siempre con el mismo
criterio.

\textbf{Configuración ANN explícita.} HNSW \cite{malkov2018} con $m{=}24$,
\code{ef\_construct}${=}120$ y \code{ef\_search}${=}128$, fijados en
\code{config/run\_config.yaml}.

\textbf{Una lección de producción: el índice que no existía.} La primera
versión funcional del sistema pasaba todas las pruebas y reportaba fidelidad
ANN 1.0\ldots{} porque no había índice. Qdrant solo construye el grafo HNSW
cuando un segmento supera \code{indexing\_threshold} (10--20 MB por
defecto), y el catálogo de entonces ($15.000 \times 384$ float32) repartido
en 4 segmentos no lo alcanzaba: estado \emph{green} con \code{indexed\_vectors\_count}${=}0$, y
cada búsqueda ``ANN'' era en realidad un escaneo exhaustivo con
\code{ef\_search} inerte. Una revisión adversarial del código lo detectó
verificándolo contra la colección viva. La corrección tiene dos partes:
\code{indexing\_threshold}${=}100$~KB al crear la colección, y una
verificación tras la ingesta ---previa a servir consultas del pipeline--- que
exige estado \emph{green} \textbf{y} una cobertura mínima del 80\,\% de
vectores indexados (esperar solo el estado verde no detecta nada, porque
verde no implica indexado; en la ejecución final la cobertura es
15.000/15.000). Tras la corrección, las latencias del
escaneo se redujeron aproximadamente a la mitad con la configuración de
entonces (384d: p50 de ${\sim}4.3$ a ${\sim}2.3$~ms). La moraleja queda en
el sistema: la verificación del estado de indexación debe comprobar el
índice, no el semáforo.

\textbf{Ingesta por lotes e idempotente.} Lotes de 256 con \code{wait=True}
y \emph{upsert} por \code{record\_id}: repetir la ingesta completa deja
exactamente 15.000 registros (verificado en el informe de ingesta:
\code{count\_before}${=}$\code{count\_after}${=}15000$,
\code{idempotent=true}). La ingesta completa tarda ${\sim}4$~s en local y la
colección puede reconstruirse desde cero con \code{AURUM\_ALLOW\_RESET=true}.

\section{Recuperación y filtros}

La búsqueda global devuelve top-$k$ con score nativo; la filtrada añade una
condición \code{brand equals <valor>} \textbf{dentro} de
\code{query\_points} (filtrado del lado del servidor, sobre el índice de
payload), nunca un post-filtrado del top-10 global. Las cuatro consultas
filtradas devuelven 10/10 resultados que cumplen la marca (Einhell, Apple,
NIKE, SAMSUNG), verificado automáticamente en el informe de filtros de
\code{resultados/evidencia/}.

Los casos límite tienen tratamiento explícito y probado: una colección vacía
lanza un error accionable (``ingiere con \code{make ingest}''), un filtro
sin resultados devuelve lista vacía sin error (sondeado con una marca
inexistente: 0 resultados), y un Qdrant caído se traduce en
\code{VectorStoreUnavailableError} con instrucciones, en todas las
operaciones (búsqueda, ingesta, lecturas, borrados).

\section{Operaciones sobre el catálogo}
\label{sec:operaciones}

Los 24 eventos de \code{eventos\_catalogo.csv} (8 actualizaciones, 8 altas,
8 bajas) se aplican en orden de \code{sequence} sobre una colección dedicada
sembrada con el catálogo íntegro, de modo que los artefactos entregados
---generados sobre el catálogo prístino--- siguen siendo reproducibles.
Ambas colecciones comparten esquema, configuración e ingesta.

\begin{itemize}
  \item \textbf{Idempotencia}: la secuencia completa se aplica \textbf{dos
  veces}; el recuento final es 15.000 en ambas pasadas ($8$ altas $-$ $8$
  bajas) y el estado final se verifica \emph{registro a registro}: los 24
  \code{record\_id} afectados se leen por ID y muestran la versión entregada
  o han dejado de existir.
  \item \textbf{Visibilidad}: para una operación de cada tipo se verifica
  lectura por ID y consulta vectorial con espera activa acotada
  (\code{wait\_until} con \emph{timeout}). Con \code{wait=True} en las
  escrituras, la actualización muestra \code{catalog\_version}${=}2$ y
  aparece en el top-3 de su propio vector al primer intento (${\sim}3$~ms);
  el alta es recuperable por ambas rutas; la baja deja de serlo. El sistema
  no cronometra proveedores: demuestra que una escritura confirmada acaba
  siendo observable y que sabe esperar o fallar con \emph{deadline} si no lo
  es.
\end{itemize}

\textbf{Seguridad de las operaciones.} Toda operación valida que la
colección empiece por el prefijo \code{aurum-market-eval} antes de construir
el cliente; recrear una colección exige \code{AURUM\_ALLOW\_RESET=true};
borrarla exige la confirmación exacta
\code{AURUM\_CONFIRM\_CLEANUP=DELETE:<colección>}. Todo desactivado por
defecto, sin credenciales en el repositorio y sin servicios cloud: el
corrector no hereda costes.

\section{Control de altas duplicadas}
\label{sec:duplicados}

\textbf{Regla.} La base vectorial genera los 5 candidatos más próximos del
alta entrante (codificada como documento, con la misma composición que el
catálogo). Se declara duplicado si el mejor candidato alcanza
$\mathit{score} \geq 0.937$; el margen sobre el segundo candidato se registra
como evidencia (umbral de margen $0.0$: la rejilla lo descarta, porque hay
duplicados reales casi empatados con una segunda copia legítima).

\textbf{Calibración solo con desarrollo.} Sobre los 14 casos etiquetados, la
rejilla explora todos los umbrales que separan scores observados. Los 7
duplicados puntúan en $[0.974, 1.000]$ y los 7 no duplicados en
$[0.870, 0.899]$: cualquier umbral de score en ese hueco logra $F_1{=}1.0$.
Se eligió un umbral \textbf{centrado en el margen de separación} (0.937,
punto medio 0.9365) en lugar del umbral máximo que devuelve la rejilla
(0.9734), criterio \emph{max-margin} que no se pega a ningún caso
concreto y generaliza mejor. Con él: precision 1.0, recall 1.0, $F_1$ 1.0 en
desarrollo, exigiendo además que el candidato señalado sea exactamente el
producto de referencia etiquetado.

\textbf{Costes asimétricos.} Un falso positivo retiene una ficha legítima en
revisión: molesta al vendedor y añade trabajo manual, pero es reversible en
minutos. Un falso negativo publica un duplicado que fragmenta reseñas y
stock y degrada el propio buscador; su coste es mayor y más silencioso. Por
eso el desempate de la calibración prima la precisión (la cola de revisión
debe ser creíble) pero el umbral elegido queda centrado, sin sacrificar
recall. En desarrollo no hubo ni un tipo de error ni el otro; el caso más
difícil (DEV-NEW-007, score 0.899, un no-duplicado muy parecido) queda a
0.038 del umbral y es el que habría que vigilar en producción.

\textbf{Evaluación ciega.} Sobre \code{altas\_evaluacion.csv} la regla
predice 7 duplicados y 7 altas nuevas; los positivos puntúan en
$[0.941, 1.000]$ con margen medio $0.08$ y los negativos en
$[0.862, 0.881]$: la misma separación limpia que en desarrollo, señal de que
el umbral no está sobreajustado.

\section{Evaluación}
\label{sec:evaluacion}

Todas las métricas se regeneran con un único comando (\code{make metrics}) y
quedan en \code{resultados/metricas\_desarrollo.json}; cada experimento
conserva configuración, métricas e IDs recuperados.

\textbf{Declaración de relevancia} (fijada antes de los experimentos y sin
cambios): nDCG@10 \cite{jarvelin2002} usa ganancias E=3, S=2, C=1, I=0;
\textbf{Recall@10 considera relevantes E$\cup$S} con el conjunto relevante
completo como denominador (con hasta 39 relevantes por consulta, el techo
alcanzable con $k{=}10$ es $<1.0$; se reporta sin maquillar el denominador);
\textbf{MRR@10 considera relevante solo E}. La tabla~\ref{tab:metricas}
resume la ejecución final.

\begin{table}[h]
  \centering\small
  \caption{Métricas de la configuración final medidas a través de Qdrant
  (ruta de producción). Latencias: 12 consultas $\times$ 30 repeticiones
  tras 5 de calentamiento, en local (WSL2, 8 vCPU); varían $\pm 0.3$~ms
  entre ejecuciones y no comparan proveedores.}
  \label{tab:metricas}
  \begin{tabular}{@{}lc@{}}
    \toprule
    Métrica & Valor \\
    \midrule
    nDCG@10 & 0.563 \\
    Recall@10 (E$\cup$S) & 0.248 \\
    MRR@10 (E) & 0.792 \\
    Latencia p50 / p95 & 2.71 / 3.43 ms \\
    Fidelidad ANN@10 (media / mín.) & 1.000 / 1.000 \\
    \bottomrule
  \end{tabular}
\end{table}

\textbf{La fidelidad es una medición, no una tautología.} Tras el incidente
de la sección~\ref{sec:indice}, la fidelidad ---fracción del top-10 exacto
que recupera el ANN--- se valida con un barrido de \code{ef\_search}
(\code{make sweep-ef}) sobre dos grafos del mismo catálogo: el entregado
($m{=}24$, optimizado) y uno deliberadamente débil ($m{=}4$,
\code{ef\_construct}${=}16$) donde el compromiso emerge
(tabla~\ref{tab:barrido}).

\begin{table}[h]
  \centering\small
  \caption{Barrido de \code{ef\_search}: fidelidad frente al oráculo exacto
  (20 consultas).}
  \label{tab:barrido}
  \begin{tabular}{@{}llcc@{}}
    \toprule
    Grafo & \code{ef} & Fid. media@10 & Fid. mín.@10 \\
    \midrule
    $m{=}24$ (entregado) & 10 & 0.885 & 0.30 \\
    $m{=}24$ (entregado) & 64 & 0.985 & 0.80 \\
    $m{=}24$ (entregado) & 128 & \textbf{1.000} & \textbf{1.000} \\
    $m{=}4$ (débil) & 10 & 0.435 & 0.10 \\
    $m{=}4$ (débil) & 128 & 0.805 & 0.30 \\
    \bottomrule
  \end{tabular}
\end{table}

Dos lecciones: el grafo entregado no pierde nada en todo el barrido (a esta
escala, con un grafo bien construido, \code{ef\_search} es margen de
seguridad, no dial crítico), y un grafo mal construido \textbf{no se arregla
buscando más}: con $m{=}4$ la fidelidad se estanca en ${\sim}0.87$ incluso a
\code{ef}${=}128$, porque las aristas que no existen no se pueden recorrer.
La calidad se decide en construcción ($m$, \code{ef\_construct}); al crecer
el catálogo este barrido es la herramienta de re-calibración.

\section{Atribución de errores}
\label{sec:atribucion}

Tres fallos representativos, con la capa responsable determinada por
evidencia. La fidelidad ANN 1.000 en la configuración final descarta el
índice en todos ellos: el vecino que devuelve HNSW es el vecino exacto
(verificado consulta a consulta en el notebook 05).

\textbf{Caso 1 --- Representación: \guillemotleft estantes sin taladro
habitacion\guillemotright{} (38249, nDCG 0.379).} El modelo base mejora la
cabeza del ranking (el primer resultado ya es un estante E; con e5-small el
MRR era 0.25), pero la cola sigue contaminada por libros (\emph{La buena
cocina sin sal}, \emph{Máscaras sin nombre}, \emph{Ángeles sin alas}): E5
ancla la negación ``sin + sustantivo'' y con títulos cortos sin marca el
parecido superficial domina. El oráculo exacto devuelve lo mismo: el vecino
exacto ya es semánticamente malo, luego el error nace en la representación,
no en el índice. Mitigación razonable: componer el texto con una frase de
categoría o un reranking léxico ligero del top-50.

\textbf{Caso 2 --- Datos y juicios: \guillemotleft disfraz halloween talla
grande hombre\guillemotright{} (33633, MRR 0.0).} El sistema devuelve
disfraces de Halloween reales (recall 0.50, el más alto de esta consulta en
toda la rejilla), pero la consulta tiene 16 juicios (solo 4 relevantes
E$\cup$S) y el único Exact etiquetado es\ldots{} una blusa de mujer (ruido
del dataset ESCI \cite{reddy2022}): el MRR es 0 precisamente porque el
sistema \emph{no} recupera esa blusa, y los disfraces plausibles sin juicio
computan ganancia 0. Es un error de la capa de datos/etiquetas que ninguna
configuración puede remontar.

\textbf{Caso 3 --- Representación en atributos finos: \guillemotleft botines
marrones mujer tacon medio\guillemotright{} (18868, nDCG 0.370).} En el
top-5 conviven un botín \emph{plateado de tacón bajo} etiquetado Irrelevant
y botines de tacón alto sin juicio: la categoría se acierta, pero color y
altura de tacón se difuminan en la similitud coseno --- y el patrón persiste
al pasar de e5-small a e5-base, señal de que no es cuestión de capacidad. La
mitigación estructural no es cambiar de modelo sino mover atributos duros
(color) a metadatos filtrables, como ya se hace con la marca.

\textbf{Las otras dos capas, con su propia evidencia.} La capa de
\textbf{índice} se demuestra provocándole una pérdida real: en el grafo
débil del barrido ($m{=}4$, \code{ef\_search}${=}10$) el ANN pierde hasta 9
de los 10 vecinos que el oráculo recupera en una consulta, exactamente la
forma de evidencia que define el enunciado; en la configuración entregada la
pérdida es cero. En \textbf{persistencia/consistencia} no se observó
discrepancia alguna: las escrituras usan \code{wait=True}, las tres
verificaciones de visibilidad de la sección~\ref{sec:operaciones} convergen
al primer intento y el estado final de los 24 eventos se valida registro a
registro.

\section{Decisión recomendada y evolución}

\textbf{Para Aurum Market hoy:} Qdrant local con HNSW ($m{=}24$,
\code{ef\_construct}${=}120$, \code{ef\_search}${=}128$),
\code{multilingual-e5-base} sobre título+marca+color, filtro de marca por
payload indexado, regla de duplicados con umbral 0.937 sobre score coseno y
revisión humana de los positivos. Todo el recorrido corre en un portátil, es
idempotente y se reconstruye desde cero con dos comandos.

\textbf{Qué cambiaría al crecer el catálogo (15k $\rightarrow$ 1M):}

\begin{itemize}
  \item \emph{Índice}: el barrido de \code{ef\_search} pasa de curiosidad a
  herramienta de operación; habría que re-medir fidelidad por percentil de
  consulta y presupuestar memoria del grafo (o evaluar cuantización escalar
  de Qdrant).
  \item \emph{Representación}: el caso 1 de la
  sección~\ref{sec:atribucion} escala mal; un modelo mayor (e5-large), la
  opción API ya preparada (Gemini Embedding 2) o un reranker sobre el top-100
  se re-evaluarían con el mismo arnés de experimentos, que ya conserva IDs y
  configuraciones.
  \item \emph{Duplicados}: con más altas diarias, el umbral único debería
  convertirse en dos (auto-rechazo / revisión humana / auto-aprobación)
  calibrados sobre la curva precision-recall completa.
  \item \emph{Operativa}: la colección de eventos separada se sustituiría
  por \emph{snapshots} de Qdrant y un flujo de reingesta azul/verde; los
  interruptores de limpieza ya existen y se mantendrían.
\end{itemize}

El objetivo no era demostrar que una base vectorial devuelve algo: era saber
cuándo sus resultados son útiles (consultas con intención clara y atributos
blandos), qué capa explica sus errores (representación y juicios, no el
índice) y qué decisión resiste producción (la configuración completa está en
\code{config/run\_config.yaml} y cada número de este informe se regenera
desde ella).

\appendix

\section{Reproducción y artefactos}

\textbf{Reproducción íntegra} (entorno limpio, ${\sim}10$ min sin GPU):
\code{make setup \&\& make up \&\& make embeddings \&\& make pipeline}.

\textbf{Notebooks de I+D (entregados ejecutados).} La serie
\code{notebooks/actividad\_00}--\code{05} justifica cada bloque de
decisiones contra el sistema vivo ---caso/datos/baseline, representación,
índice y BD, recuperación y filtros, operaciones y duplicados, y evaluación
completa con atribución de errores y el checklist del enunciado verificado
por aserciones--- y explica cada concepto en su primera aparición para
lectura secuencial. Se regeneran con \code{make notebook} y se ejecutan
\emph{headless} con \code{make execute-notebook}.

\textbf{Artefactos entregados.}

\begin{description}[style=nextline,itemsep=2pt,font=\normalfont]
  \item[\code{resultados/resultados\_busqueda.csv}] Top-10 por consulta ciega
  ($12 \times 10$ filas, IDs únicos por consulta).
  \item[\code{resultados/resultados\_duplicados.csv}] Decisión, candidato y
  score de las 14 altas ciegas.
  \item[\code{resultados/metricas\_desarrollo.json}] nDCG/Recall/MRR@10,
  latencias p50/p95, fidelidad ANN, duplicados y entorno.
  \item[\code{config/run\_config.yaml}] Configuración exacta de la ejecución
  final.
  \item[\code{docs/images/arquitectura.svg}] Diagrama de arquitectura
  (figura~\ref{fig:arquitectura}).
  \item[\code{resultados/evidencia/}] Copia versionada de la evidencia:
  registro de experimentos con IDs recuperados, evaluación completa con
  entorno, barrido \code{ef\_search}, calibración y decisiones de duplicados,
  informes de filtros, eventos e ingesta.
  \item[\code{.artifacts/}] Los mismos informes en su ubicación de trabajo
  (regenerables con \code{make pipeline}).
\end{description}

\textbf{Verificación previa a la entrega.} Ingesta repetida sin aumentar
recuento \checkmark{} · filtros nunca devuelven otra marca \checkmark{} ·
eventos dejan el estado esperado dos veces \checkmark{} · 12 rankings ciegos
con 10 IDs únicos y válidos \checkmark{} · toda predicción positiva señala
candidato \checkmark{} · métricas regenerables con \code{make metrics}
\checkmark{} · sin claves ni datos reservados en el repositorio \checkmark{}
(55 pruebas automatizadas: \code{make test}, \code{make test-integration}).

\balance
\begin{thebibliography}{9}\small
\bibitem{reddy2022} C.~K. Reddy \emph{et al.} (2022). Shopping Queries
Dataset: A Large-Scale ESCI Benchmark for Improving Product Search.
\emph{arXiv:2206.06588}.
\bibitem{wang2024} L.~Wang \emph{et al.} (2024). Multilingual E5 Text
Embeddings: A Technical Report. \emph{arXiv:2402.05672}.
\bibitem{reimers2019} N.~Reimers y I.~Gurevych (2019). Sentence-BERT:
Sentence Embeddings using Siamese BERT-Networks. \emph{arXiv:1908.10084}.
\bibitem{malkov2018} Y.~A. Malkov y D.~A. Yashunin (2018). Efficient and
robust approximate nearest neighbor search using Hierarchical Navigable
Small World graphs. \emph{IEEE TPAMI}, 42(4). \emph{arXiv:1603.09320}.
\bibitem{robertson2009} S.~Robertson y H.~Zaragoza (2009). The
Probabilistic Relevance Framework: BM25 and Beyond. \emph{Foundations and
Trends in Information Retrieval}, 3(4).
\bibitem{jarvelin2002} K.~Järvelin y J.~Kekäläinen (2002). Cumulated
gain-based evaluation of IR techniques. \emph{ACM TOIS}, 20(4).
\bibitem{qdrant} Qdrant. Documentación: colecciones, HNSW y filtrado por
payload. \url{https://qdrant.tech/documentation/concepts/}.
\end{thebibliography}

\end{document}
